\section{Aggregation Bias Across COICOP Digits}
\label{app:aggregation-bias}

This appendix assesses how the number of COICOP digits used to classify products affects the measured inflation gap
between households in the first and fifth income quintiles. In contrast to a calculation based on
country-level group baskets, Table~\ref{tab:individual-distribution-aggregation-bias-2022-excluding-rents}
is constructed from the household-level HBS microdata used in the individual inflation-distribution
pipeline and reports the pipeline's pooled income-quintile specification. Germany, France, Spain,
the Netherlands, and Belgium are included. Italy is excluded
because income is unavailable in the HBS file and households therefore cannot be assigned to income
quintiles.

\paragraph{Construction of the table.}
Households are ranked by equivalised disposable income, defined as reported household net income
divided by the modified-OECD equivalence scale. Survey-weighted quintile thresholds are calculated
on the pooled five-country HBS sample, as in the individual-distribution pipeline. The first and
fifth quintiles are denoted by $Q1$ and $Q5$. As a first step, actual rents (COICOP 041) and imputed
rents (COICOP 042) are removed from both the household expenditure universe and the price universe.
The remaining expenditures are then renormalised to sum to one within each household.

For household $h$ and COICOP category $j$ at digit depth $d$, the fixed 2020 budget share is
\begin{equation}
 s^{d}_{hj}
 =
 \frac{x^{d}_{hj,2020}}
 {\sum_{k \in \mathcal{J}^{d}_{h}}x^{d}_{hk,2020}},
\end{equation}
where $\mathcal{J}^{d}_{h}$ contains the non-rent categories with matched price observations. All
coarser baskets are reconstructed from the same finest set of expenditure variables available for a
given country. Consequently, changes across digit depths reflect product aggregation rather than a change
in the underlying HBS expenditure universe.

For each product, annual inflation in 2022 is the arithmetic mean of the available monthly
year-on-year rates:
\begin{equation}
 \bar{\pi}^{d}_{j,2022}
 =
 \frac{1}{M_j}\sum_{m=1}^{M_j}
 \left(\frac{P^{d}_{j,2022,m}}{P^{d}_{j,2021,m}}-1\right)100,
\end{equation}
with at least ten matched months required. Household-specific inflation is then
\begin{equation}
 \pi^{d}_{h,2022}
 =
 \sum_{j \in \mathcal{J}^{d}_{h}}s^{d}_{hj}\bar{\pi}^{d}_{j,2022}.
\end{equation}
The inflation rate of quintile $q$ is the HBS survey-weighted mean of household inflation,
$\pi^{d}_{q,2022}$. The reported gap and aggregation bias are
\begin{align}
 G^{d}_{2022} &= \pi^{d}_{Q1,2022}-\pi^{d}_{Q5,2022},\\
 B^{d,d+1}_{2022} &= G^{d}_{2022}-G^{d+1}_{2022}.
\end{align}
A positive gap means that low-income households experienced higher inflation than high-income
households. A negative aggregation bias means that moving to the finer classification increases the
measured $Q1$--$Q5$ gap.


\begin{longtable}[t]{llrrrrr}
\caption{\label{tab:individual-distribution-aggregation-bias-2022-excluding-rents}Aggregation bias excluding actual and imputed rents, 2022}\\
\toprule
Definition & Country & Digit 2 gap & Digit 3 gap & Digit 4 gap & D2--D3 bias & D2--D4 bias\\
\midrule
Pooled pipeline quintiles & DE & 1.2 & 1.9 & -- & -0.6 & --\\
Pooled pipeline quintiles & FR & 0.0 & 0.2 & -- & -0.2 & --\\
Pooled pipeline quintiles & ES & 0.3 & 0.9 & -- & -0.6 & --\\
Pooled pipeline quintiles & NL & -0.2 & 4.1 & 5.4 & -4.2 & -5.5\\
Pooled pipeline quintiles & BE & 0.6 & 1.5 & 2.2 & -0.9 & -1.6\\
\bottomrule
\end{longtable}
\begin{flushleft}
\footnotesize{\emph{Notes:} Household inflation is calculated using each household's expenditure shares and the mean of monthly year-on-year item inflation rates in 2022, matching the individual-distribution pipeline. Q1 and Q5 means use HBS survey weights and reproduce the pooled income-quintile convention in the distribution script. Actual rents (041) and imputed rents (042) are excluded before household shares are normalised. COICOP digits 2--3 are available for DE, FR and ES, and digits 2--4 for NL and BE.}
\end{flushleft}


The finer classification increases the estimated $Q1$--$Q5$ gap in all five countries. For the
Netherlands and Belgium, moving from digit 3 to digit 4 raises it by a further 1.315 and 0.674
percentage points. Digit 4 is unavailable for Germany, France, and Spain because their harmonised
HBS files stop at digit 3. Since product-price matching declines at finer digit depths, the estimates
combine aggregation with residual COICOP coverage limits. They are a measurement diagnostic, not a
measure of household substitution.
